% !TEX program = pdflatex
\documentclass[12pt,a4paper]{article}

% Pacotes essenciais
\usepackage[utf8]{inputenc}
\usepackage[T1]{fontenc}
\usepackage[brazil]{babel}
\usepackage{amsmath,amssymb,amsthm}
\usepackage{graphicx}
\usepackage{hyperref}
\usepackage{geometry}
\usepackage{booktabs}
\usepackage{array}
\usepackage{multirow}
\usepackage{longtable}
\usepackage{algorithm}
\usepackage{algorithmic}
\usepackage{listings}
\usepackage{xcolor}
\usepackage{tcolorbox}
\usepackage{enumitem}
\usepackage{fancyhdr}
\usepackage{float}

% Configurações de página
\geometry{
    a4paper,
    left=2.5cm,
    right=2.5cm,
    top=3cm,
    bottom=3cm
}

% Configuração de cabeçalho e rodapé
\pagestyle{fancy}
\fancyhf{}
\fancyhead[L]{Análise Comparativa: Sequencial vs PL v2.0}
\fancyhead[R]{Sistema SIVIRA}
\fancyfoot[C]{\thepage}

% Configuração de hyperlinks
\hypersetup{
    colorlinks=true,
    linkcolor=blue,
    filecolor=magenta,
    urlcolor=cyan,
    citecolor=blue,
    pdftitle={Análise Comparativa Sequencial vs PL},
    pdfauthor={Sistema SIVIRA},
}

% Configuração de código
\lstset{
    basicstyle=\ttfamily\small,
    keywordstyle=\color{blue},
    commentstyle=\color{gray},
    stringstyle=\color{red},
    numbers=left,
    numberstyle=\tiny\color{gray},
    stepnumber=1,
    numbersep=5pt,
    backgroundcolor=\color{white},
    showspaces=false,
    showstringspaces=false,
    showtabs=false,
    frame=single,
    tabsize=2,
    captionpos=b,
    breaklines=true,
    breakatwhitespace=false,
}

% Definições de teoremas
\theoremstyle{definition}
\newtheorem{definition}{Definição}[section]
\newtheorem{theorem}{Teorema}[section]
\newtheorem{lemma}[theorem]{Lema}
\newtheorem{proposition}[theorem]{Proposição}

% Comandos customizados
\newcommand{\BigO}{\mathcal{O}}
\newcommand{\reals}{\mathbb{R}}
\newcommand{\integers}{\mathbb{Z}}
\newcommand{\naturals}{\mathbb{N}}

% Título
\title{
    \textbf{Análise Comparativa:}\\
    \textbf{Método Sequencial vs Programação Linear v2.0}\\
    \vspace{0.5cm}
    \large Sistema SIVIRA - Scheduling de Produção em Padaria Industrial
}

\author{
    Sistema SIVIRA\\
    \texttt{www.sivira.com}
}

\date{14 de Novembro de 2025}

\begin{document}

\maketitle

\begin{abstract}
Este documento apresenta uma análise comparativa rigorosa entre dois métodos de \textit{scheduling} de produção implementados no sistema SIVIRA: (1) Método Sequencial baseado em heurística \textit{greedy} com \textit{backward scheduling} e ordenação EDD; e (2) Programação Linear v2.0 utilizando Mixed Integer Programming (MIP) com solver SCIP. Os resultados empíricos com 13 pedidos reais mostram que o método sequencial atinge 84.6\% de taxa de sucesso (11/13 pedidos) com \textit{makespan} de 29h em 0.42s, enquanto o PL v2.0 atinge 69.2\% (9/13 pedidos) com \textit{makespan} de 73h em 0.18s. Apesar da garantia de otimalidade do PL, o método sequencial apresenta melhor desempenho prático devido ao espaço de busca mais amplo e função objetivo implicitamente multi-critério. O documento discute este paradoxo, analisa trade-offs fundamentais e propõe melhorias para ambos os métodos.
\end{abstract}

\tableofcontents

\newpage

\section{Introdução e Contexto}

\subsection{Problema de Scheduling de Produção}

O sistema SIVIRA resolve um problema de \textbf{Job Shop Scheduling} com características específicas de uma padaria industrial. Este problema pode ser formalmente classificado como:

\begin{equation}
Jm \mid prec, r_j, d_j \mid C_{max} \text{ ou } \sum U_j
\end{equation}

onde:
\begin{itemize}
    \item $Jm$: Job Shop com $m$ máquinas
    \item $prec$: Restrições de precedência entre atividades
    \item $r_j$: Datas de liberação (release dates)
    \item $d_j$: Prazos (due dates)
    \item $C_{max}$: Makespan (minimizar)
    \item $\sum U_j$: Número de trabalhos atrasados (minimizar)
\end{itemize}

Este é um problema \textbf{NP-difícil} \cite{Pinedo2016}, o que justifica o uso de heurísticas eficientes além de métodos exatos.

\subsection{Características do Problema}

O problema apresenta as seguintes características:

\begin{enumerate}
    \item \textbf{Múltiplos pedidos} com diferentes produtos
    \item \textbf{Atividades sequenciais} com precedências estritas
    \item \textbf{Recursos compartilhados} (equipamentos com capacidade limitada)
    \item \textbf{Restrições temporais}:
    \begin{itemize}
        \item Deadlines por pedido
        \item Tempo máximo de espera entre atividades
        \item Janelas de produção (horário comercial)
    \end{itemize}
    \item \textbf{Objetivo}: Maximizar número de pedidos atendidos
\end{enumerate}

\subsection{Abordagens Implementadas}

Duas abordagens foram desenvolvidas e comparadas:

\begin{description}
    \item[Método Sequencial] Heurística \textit{greedy} com complexidade $\BigO(n \log n + n^2)$
    \item[Programação Linear v2.0] Modelo MIP com complexidade NP-completo
\end{description}

\section{Formulação Matemática do PL v2.0}

\subsection{Notação e Conjuntos}

Definimos os seguintes conjuntos e parâmetros:

\begin{table}[H]
\centering
\begin{tabular}{@{}lll@{}}
\toprule
\textbf{Símbolo} & \textbf{Descrição} & \textbf{Cardinalidade} \\
\midrule
$P$ & Conjunto de pedidos & $|P| = n$ \\
$J_p$ & Janelas viáveis para pedido $p \in P$ & $|J_p| \leq j_{max}$ \\
$A_p$ & Atividades do pedido $p$ & $|A_p| = m_p$ \\
$E$ & Conjunto de equipamentos & $|E| = e$ \\
$T$ & Conjunto de slots temporais & $|T| = t$ \\
\bottomrule
\end{tabular}
\caption{Conjuntos básicos do modelo}
\label{tab:conjuntos}
\end{table}

\textbf{Parâmetros de configuração:}
\begin{itemize}
    \item $j_{max} = 15$: Número máximo de janelas por pedido
    \item $\Delta t = 30$: Resolução temporal em minutos
    \item $t_{timeout} = 600$: Timeout do solver em segundos
\end{itemize}

\subsection{Variáveis de Decisão}

A variável principal do modelo é binária:

\begin{equation}
x_{p,j} \in \{0, 1\} \quad \forall p \in P, j \in J_p
\end{equation}

\textbf{Interpretação:}
\begin{equation}
x_{p,j} = \begin{cases}
1 & \text{se pedido } p \text{ usa janela } j \\
0 & \text{caso contrário}
\end{cases}
\end{equation}

O número total de variáveis é:
\begin{equation}
|X| = \sum_{p \in P} |J_p|
\end{equation}

\textbf{Exemplo prático:} Para $n=13$ pedidos com $|J_p|=9$ janelas cada, temos $|X|=117$ variáveis.

\subsection{Função Objetivo}

O objetivo é maximizar o número de pedidos atendidos:

\begin{equation}
\max \sum_{p \in P} y_p
\end{equation}

onde $y_p$ indica se o pedido $p$ foi atendido:

\begin{equation}
y_p = \begin{cases}
1 & \text{se } \exists j \in J_p : x_{p,j} = 1 \\
0 & \text{caso contrário}
\end{cases}
\end{equation}

Como a restrição de unicidade (Equação \ref{eq:unicidade}) garante $\sum_j x_{p,j} \leq 1$, podemos linearizar:

\begin{equation}
\max \sum_{p \in P} \sum_{j \in J_p} x_{p,j}
\label{eq:objetivo}
\end{equation}

\subsection{Restrições}

\subsubsection{Restrição de Unicidade de Janela}

Cada pedido pode usar no máximo uma janela temporal:

\begin{equation}
\sum_{j \in J_p} x_{p,j} \leq 1 \quad \forall p \in P
\label{eq:unicidade}
\end{equation}

Esta restrição gera $|P| = n$ inequações lineares.

\subsubsection{Restrições de Tempo Máximo de Espera}

Para pedidos com atividades sucessoras que têm tempo máximo de espera $\tau_{max}(a_i)$:

\begin{equation}
gap(a_i, a_{i+1}) \leq \tau_{max}(a_i) \quad \forall i \in \{1, ..., m-1\}
\label{eq:gap}
\end{equation}

onde:
\begin{equation}
gap(a_i, a_{i+1}) = inicio(a_{i+1}) - fim(a_i)
\end{equation}

\textbf{Caso especial} ($\tau_{max} = 0$, atividades contíguas):
\begin{equation}
inicio(a_{i+1}) = fim(a_i)
\end{equation}

\subsubsection{Restrições de Conflitos Temporais}

Duas janelas $j_1$ e $j_2$ se \textbf{sobrepõem} se:

\begin{equation}
overlap(j_1, j_2) \iff \begin{cases}
inicio_{j_1} < fim_{j_2} \\
inicio_{j_2} < fim_{j_1}
\end{cases}
\label{eq:overlap}
\end{equation}

Para cada par de janelas que se sobrepõem:

\begin{equation}
x_{p_1,j_1} + x_{p_2,j_2} \leq 1 \quad \forall p_1 \neq p_2, j_1 \in J_{p_1}, j_2 \in J_{p_2} : overlap(j_1, j_2)
\label{eq:conflito}
\end{equation}

O número de restrições de conflito no pior caso é:

\begin{equation}
R_{conflitos} = \binom{|J_{total}|}{2} = \frac{|J_{total}| \times (|J_{total}| - 1)}{2}
\end{equation}

\subsection{Modelo Completo}

O modelo completo de programação linear inteira é:

\begin{equation}
\begin{aligned}
\max \quad & \sum_{p \in P} \sum_{j \in J_p} x_{p,j} \\
\text{s.a.} \quad & \sum_{j \in J_p} x_{p,j} \leq 1 && \forall p \in P \\
& x_{p_1,j_1} + x_{p_2,j_2} \leq 1 && \forall (p_1,j_1), (p_2,j_2) : overlap(j_1,j_2) \\
& x_{p,j} \in \{0, 1\} && \forall p \in P, j \in J_p
\end{aligned}
\label{eq:modelo_completo}
\end{equation}

\subsection{Estatísticas do Modelo}

Para o conjunto de dados com 13 pedidos, o modelo apresenta as seguintes características:

\begin{table}[H]
\centering
\begin{tabular}{@{}lcc@{}}
\toprule
\textbf{Métrica} & \textbf{Otimizado} & \textbf{Antigo} \\
\midrule
Variáveis & 117 & 65 \\
Restrições Totais & 1,321 & 513 \\
$\quad$ Unicidade & 13 & 13 \\
$\quad$ Tempo Máx Espera & 0 & 0 \\
$\quad$ Equipamentos & 0 & 0 \\
$\quad$ Conflitos Temporais & 1,308 & 500 \\
\midrule
Status Solver & OPTIMAL & OPTIMAL \\
Tempo Resolução & 0.18s & 0.04s \\
\bottomrule
\end{tabular}
\caption{Estatísticas do modelo PL v2.0}
\label{tab:estatisticas_modelo}
\end{table}

\section{Método Sequencial (Backward Scheduling)}

\subsection{Descrição do Algoritmo}

O método sequencial é uma heurística \textit{greedy} baseada em três etapas:

\begin{enumerate}
    \item \textbf{Ordenação} dos pedidos por deadline (EDD - Earliest Due Date)
    \item \textbf{Backward scheduling} para cada pedido
    \item \textbf{Alocação sequencial} de equipamentos
\end{enumerate}

\subsection{Pseudo-código}

\begin{algorithm}[H]
\caption{Scheduling Sequencial}
\label{alg:sequencial}
\begin{algorithmic}[1]
\REQUIRE Lista de pedidos $P = \{p_1, p_2, ..., p_n\}$
\ENSURE Schedule $S$ ou FALHA
\STATE $P_{ordenado} \leftarrow \text{SORT}(P, key=\lambda p: p.deadline)$
\FOR{cada pedido $p$ em $P_{ordenado}$}
    \STATE $A_p \leftarrow \text{CRIAR\_ATIVIDADES}(p)$
    \STATE $inicio_{corrente} \leftarrow p.deadline$
    \FOR{cada atividade $a$ em $\text{REVERSO}(A_p)$}
        \STATE $fim_{desejado} \leftarrow inicio_{corrente}$
        \STATE $inicio_{desejado} \leftarrow fim_{desejado} - a.duracao$
        \STATE $sucesso \leftarrow \text{ALOCAR\_EQUIPAMENTOS}(a, inicio_{desejado}, fim_{desejado})$
        \IF{não $sucesso$}
            \STATE \textbf{rejeitar} pedido $p$
            \STATE \textbf{continuar} com próximo pedido
        \ENDIF
        \STATE $inicio_{corrente} \leftarrow inicio_{real}$ de $a$
    \ENDFOR
    \STATE $\text{REGISTRAR\_SUCESSO}(p)$
\ENDFOR
\RETURN schedule $S$
\end{algorithmic}
\end{algorithm}

\subsection{Backward Scheduling Formal}

Para um pedido com atividades $\{a_1, a_2, ..., a_m\}$ e deadline $d$:

\begin{align}
fim(a_m) &\leq d \\
inicio(a_m) &= fim(a_m) - duracao(a_m) \\
fim(a_{m-1}) &= inicio(a_m) - gap(a_{m-1}, a_m) \\
inicio(a_{m-1}) &= fim(a_{m-1}) - duracao(a_{m-1}) \\
&\vdots \nonumber \\
inicio(a_1) &= fim(a_1) - duracao(a_1)
\end{align}

\subsection{Análise de Complexidade}

\subsubsection{Complexidade Temporal}

\begin{enumerate}
    \item \textbf{Ordenação EDD}: $\BigO(n \log n)$
    \item \textbf{Para cada pedido}:
    \begin{itemize}
        \item Criar atividades: $\BigO(m)$
        \item Backward scheduling: $\BigO(m)$
        \item Para cada atividade:
        \begin{itemize}
            \item Buscar equipamento: $\BigO(e)$
            \item Buscar slot: $\BigO(s)$
            \item Alocar: $\BigO(1)$
        \end{itemize}
    \end{itemize}
\end{enumerate}

Complexidade total:
\begin{equation}
T_{seq}(n) = \BigO(n \log n) + \sum_{i=1}^{n} \BigO(m_i \times (e + s_i))
\end{equation}

No pior caso (todos executados):
\begin{equation}
T_{seq}(n) = \BigO(n \log n + n^2)
\end{equation}

\textbf{Classe de complexidade:} Polinomial $\BigO(n^2)$

\subsubsection{Complexidade Espacial}

\begin{equation}
S_{seq}(n) = \BigO(n \times m \times e) \approx \BigO(n)
\end{equation}

assumindo $m$ e $e$ constantes.

\section{Comparação Empírica}

\subsection{Conjunto de Dados}

O dataset consiste em 13 pedidos reais de uma padaria industrial, cobrindo o período de 28/12 07:00 até 31/12 08:00 (72 horas). Os pedidos incluem pães com fermentação longa (19-27h) e salgados de produção rápida (2-5h).

\subsection{Resultados Quantitativos}

\begin{table}[H]
\centering
\begin{tabular}{@{}lccc@{}}
\toprule
\textbf{Métrica} & \textbf{Sequencial} & \textbf{PL v2.0} & \textbf{Diferença} \\
\midrule
Pedidos Atendidos & 11/13 & 9/13 & -2 pedidos \\
Taxa de Sucesso & \textbf{84.6\%} & 69.2\% & -15.4 p.p. \\
Tempo Resolução & 0.42s & \textbf{0.18s} & -57\% \\
Makespan & \textbf{29h} & 73h & +151\% \\
Pedidos Rejeitados & P2, P3 & P2, P3, P11, P13 & +2 \\
Complexidade & N/A & 117 vars, 1321 const & -- \\
Status & Completo & \textbf{OPTIMAL} & -- \\
Garantias & Nenhuma & Otimalidade & \checkmark \\
\bottomrule
\end{tabular}
\caption{Comparação quantitativa dos métodos}
\label{tab:comparacao_resultados}
\end{table}

\subsection{Análise de Makespan}

O método sequencial produz um \textit{schedule} 2.5× mais compacto (29h vs 73h), resultando em melhor utilização de recursos:

\begin{table}[H]
\centering
\begin{tabular}{@{}lccc@{}}
\toprule
\textbf{Equipamento} & \textbf{Sequencial} & \textbf{PL v2.0} & \textbf{Diferença} \\
\midrule
Forno 1 & 78\% & 52\% & -26 p.p. \\
Misturador 1 & 85\% & 48\% & -37 p.p. \\
Modeladora 1 & 92\% & 61\% & -31 p.p. \\
Fermentador 1 & 100\% & 73\% & -27 p.p. \\
\bottomrule
\end{tabular}
\caption{Utilização de recursos}
\label{tab:utilizacao_recursos}
\end{table}

\subsection{Análise de Tempo Computacional}

\begin{table}[H]
\centering
\begin{tabular}{@{}lcc@{}}
\toprule
\textbf{Fase} & \textbf{Sequencial} & \textbf{PL v2.0} \\
\midrule
Setup & 0.15s & 0.09s \\
Core (resolução) & 0.25s & \textbf{0.07s} \\
Pós-processamento & 0.02s & 0.02s \\
\midrule
\textbf{Total} & 0.42s & \textbf{0.18s} \\
\bottomrule
\end{tabular}
\caption{Perfil de tempo de execução}
\label{tab:tempo_execucao}
\end{table}

\section{Discussão e Trade-offs}

\subsection{Paradoxo da Otimalidade}

Observamos um resultado contra-intuitivo: o método PL v2.0, apesar de matematicamente ótimo, apresentou \textbf{performance inferior} ao sequencial (69.2\% vs 84.6\%).

\begin{tcolorbox}[colback=yellow!10!white,colframe=orange!75!black,title=Explicação do Paradoxo]
\begin{enumerate}
    \item \textbf{Otimalidade relativa}: PL é ótimo apenas \textit{dentro do espaço de soluções} definido pelas janelas geradas
    \item \textbf{Função objetivo limitada}: Maximizar número de pedidos não considera makespan
    \item \textbf{Discretização temporal}: Resolução de 30min pode perder soluções contínuas
    \item \textbf{Exploração do espaço}: Sequencial testa todas as posições possíveis (busca em resolução de 1min)
\end{enumerate}
\end{tcolorbox}

\subsection{Trade-offs Fundamentais}

Identificamos três trade-offs principais:

\subsubsection{Otimalidade vs Eficiência}

\begin{itemize}
    \item \textbf{Sequencial}: Eficiente ($\BigO(n^2)$) mas sem garantias
    \item \textbf{PL v2.0}: Garantias formais mas NP-completo
    \item \textbf{Gap}: Espaço para métodos híbridos
\end{itemize}

\subsubsection{Makespan vs Número de Pedidos}

Existe um trade-off fundamental não capturado pela função objetivo atual:
\begin{itemize}
    \item \textbf{Sequencial}: 11 pedidos em 29h (makespan compacto)
    \item \textbf{PL v2.0}: 9 pedidos em 73h (makespan esparso)
\end{itemize}

\subsubsection{Precisão Temporal vs Complexidade}

\begin{table}[H]
\centering
\small
\begin{tabular}{@{}lcccc@{}}
\toprule
\textbf{Resolução} & \textbf{Variáveis} & \textbf{Tempo} & \textbf{Qualidade} \\
\midrule
60 min & 65 & 0.04s & 38.5\% (5/13) \\
\textbf{30 min} & \textbf{117} & \textbf{0.18s} & \textbf{69.2\% (9/13)} \\
15 min (proj) & $\sim$210 & $\sim$2.5s & $\sim$75\% \\
5 min (proj) & $\sim$630 & $\sim$45s & $\sim$80\% \\
\bottomrule
\end{tabular}
\caption{Trade-off resolução temporal vs complexidade}
\label{tab:tradeoff_resolucao}
\end{table}

\section{Propostas de Melhoria}

\subsection{Melhorias para PL v2.0}

\subsubsection{Função Objetivo Multi-critério}

Proposta de nova função objetivo:

\begin{equation}
\max \left( w_1 \cdot \sum_p y_p - w_2 \cdot C_{max} - w_3 \cdot \sum_p L_p \right)
\end{equation}

onde:
\begin{itemize}
    \item $w_1 = 100$: Peso para número de pedidos
    \item $w_2 = 0.1$: Peso para makespan
    \item $w_3 = 10$: Peso para atrasos ($L_p = \max(0, C_p - d_p)$)
\end{itemize}

\textbf{Impacto esperado}: Reduzir makespan de 73h para $\sim$35h mantendo ou melhorando pedidos atendidos.

\subsubsection{Geração Adaptativa de Janelas}

Estratégia diferenciada por tipo de pedido:

\begin{equation}
pontos = \begin{cases}
[0.0, 0.05, 0.1, 0.15, 0.2] & \text{se } duracao_p > 20h \\
[0.0, 0.125, ..., 0.875, 1.0] & \text{caso contrário}
\end{cases}
\end{equation}

\textbf{Impacto esperado}: Melhorar atendimento de pedidos longos, aumentando taxa de 69\% para $\sim$80\%.

\subsubsection{Método Híbrido (Warm Start)}

Algoritmo proposto:

\begin{algorithm}[H]
\caption{Método Híbrido}
\begin{algorithmic}[1]
\STATE $solucao_{seq} \leftarrow$ SequencialRapido()
\STATE $warm\_start \leftarrow$ ConverterParaPL($solucao_{seq}$)
\STATE $solucao_{pl} \leftarrow$ PL\_v2.0($warm\_start$)
\RETURN $\max(solucao_{seq}, solucao_{pl})$
\end{algorithmic}
\end{algorithm}

\textbf{Impacto esperado}: Garantir pelo menos performance do sequencial, com possibilidade de melhorias via PL.

\subsection{Melhorias para Método Sequencial}

\subsubsection{Heurísticas de Ordenação Alternativas}

Além de EDD, testar:

\begin{itemize}
    \item \textbf{WSPT} (Weighted Shortest Processing Time):
    \begin{equation}
    priority_p = \frac{w_p}{p_p}
    \end{equation}

    \item \textbf{ATC} (Apparent Tardiness Cost):
    \begin{equation}
    priority_p = \frac{w_p}{p_p} \exp\left(-\frac{\max(d_p - t - p_p, 0)}{kp_{avg}}\right)
    \end{equation}
\end{itemize}

\subsubsection{Look-ahead Local}

Antes de alocar pedido $p$:
\begin{enumerate}
    \item Simular impacto em próximos 3-5 pedidos
    \item Se conflito severo detectado, testar ordem alternativa
    \item Usar busca local limitada
\end{enumerate}

\section{Conclusões e Recomendações}

\subsection{Principais Descobertas}

\begin{enumerate}
    \item Ambos os métodos são viáveis para produção real
    \item Método "ótimo" (PL) teve performance inferior ao heurístico devido a espaço de busca restrito
    \item PL v2.0 é 57\% mais rápido (0.18s vs 0.42s)
    \item Sequencial produz makespan 2.5× mais compacto
    \item Existe complementaridade entre os métodos
\end{enumerate}

\subsection{Recomendações de Uso}

\begin{table}[H]
\centering
\small
\begin{tabular}{@{}lll@{}}
\toprule
\textbf{Cenário} & \textbf{Método} & \textbf{Justificativa} \\
\midrule
$n \leq 30$ pedidos & PL v2.0 & Rápido + otimalidade \\
$n > 50$ pedidos & Sequencial & PL não escala \\
Deadline apertado & Sequencial & Makespan compacto \\
Pedidos longos (>20h) & Sequencial & PL rejeita \\
Tempo real (<1s) & PL v2.0 & Mais rápido \\
Pesquisa/análise & Ambos & Comparar \\
\bottomrule
\end{tabular}
\caption{Recomendações de uso por cenário}
\label{tab:recomendacoes}
\end{table}

\subsection{Trabalhos Futuros}

\subsubsection{Curto Prazo}
\begin{itemize}
    \item Implementar função objetivo multi-critério
    \item Desenvolver geração adaptativa de janelas
    \item Criar método híbrido (Seq + PL)
\end{itemize}

\subsubsection{Médio Prazo}
\begin{itemize}
    \item Benchmark extensivo (50+ instâncias)
    \item Análise de sensibilidade dos parâmetros
    \item Decomposição temporal para escalabilidade
\end{itemize}

\subsubsection{Longo Prazo}
\begin{itemize}
    \item Aprendizado de máquina para ordenação
    \item Otimização multi-objetivo (Pareto)
    \item Scheduling dinâmico e re-scheduling
\end{itemize}

\subsection{Mensagem Final}

\begin{tcolorbox}[colback=blue!5!white,colframe=blue!75!black,title=Conclusão Principal]
\textbf{``A escolha entre métodos depende do contexto. Não existe 'melhor método absoluto'.''}

O valor real está em \textbf{entender os trade-offs} e \textbf{escolher conscientemente} baseado nas necessidades específicas de cada situação. Completude de modelagem não garante superioridade prática.
\end{tcolorbox}

\begin{thebibliography}{9}

\bibitem{Pinedo2016}
Pinedo, M. L. (2016).
\textit{Scheduling: Theory, Algorithms, and Systems} (5th ed.).
Springer.

\bibitem{Brucker2007}
Brucker, P. (2007).
\textit{Scheduling Algorithms} (5th ed.).
Springer.

\bibitem{Wolsey1999}
Wolsey, L. A., \& Nemhauser, G. L. (1999).
\textit{Integer and Combinatorial Optimization}.
Wiley.

\bibitem{Bertsimas2005}
Bertsimas, D., \& Weismantel, R. (2005).
\textit{Optimization over Integers}.
Dynamic Ideas.

\bibitem{ORTools}
Google OR-Tools (2024).
\textit{Google Optimization Tools}.
\url{https://developers.google.com/optimization}

\end{thebibliography}

\appendix

\section{Glossário}

\begin{description}
    \item[Job Shop Scheduling] Problema de alocar trabalhos (jobs) a máquinas ao longo do tempo
    \item[Makespan] Tempo total do schedule (do início do primeiro ao fim do último trabalho)
    \item[Deadline] Prazo máximo para completar um trabalho
    \item[EDD] Earliest Due Date - Heurística que ordena por deadline crescente
    \item[Backward Scheduling] Agendar de trás para frente (do deadline para o início)
    \item[Janela Temporal] Intervalo $[inicio, fim]$ onde um trabalho pode ser executado
    \item[MIP] Mixed Integer Programming - Programação inteira mista
    \item[Branch-and-bound] Algoritmo de busca em árvore para resolver MIP
    \item[SCIP] Solving Constraint Integer Programs - Solver open-source
\end{description}

\section{Dados do Experimento}

\begin{table}[H]
\centering
\small
\begin{tabular}{@{}llccc@{}}
\toprule
\textbf{ID} & \textbf{Produto} & \textbf{Qtd} & \textbf{Deadline} & \textbf{Duração} \\
\midrule
P1 & Pão Francês & 420 & 31/12 07:00 & $\sim$6h \\
P2 & Pão Hambúrguer & 120 & 31/12 07:00 & $\sim$27h \\
P3 & Pão de Forma & 20 & 31/12 07:00 & $\sim$19h \\
P4 & Pão Baguete & 20 & 31/12 07:00 & $\sim$4h \\
P5 & Pão Trança & 15 & 31/12 07:00 & $\sim$4h \\
P6 & Coxinha Frango & 15 & 31/12 08:00 & $\sim$3h \\
P7 & Coxinha C. Sol & 10 & 31/12 08:00 & $\sim$2h \\
P8 & Coxinha Camarão & 12 & 31/12 08:00 & $\sim$3h \\
P9 & Coxinha Queijos & 12 & 31/12 08:00 & $\sim$3h \\
P10 & Folhado Frango & 10 & 31/12 08:00 & $\sim$4h \\
P11 & Folhado C. Sol & 10 & 31/12 08:00 & $\sim$3h \\
P12 & Folhado Camarão & 10 & 31/12 08:00 & $\sim$5h \\
P13 & Folhado Queijos & 5 & 31/12 08:00 & $\sim$3h \\
\bottomrule
\end{tabular}
\caption{Dataset completo do experimento}
\label{tab:dataset}
\end{table}

\end{document}
